\documentclass[10pt,a4paper]{amsart}
\usepackage[margin=19mm]{geometry}
\usepackage{amsmath,amssymb,mathtools}
\usepackage[scr=rsfs,cal=euler]{mathalfa}
\usepackage{xfrac}
\usepackage[unicode,hidelinks,bookmarks=false]{hyperref}
\newcommand{\ie}{i.e.\ }
\newcommand{\cf}{cf.\ }
\newcommand{\resp}{respectively\ }
\newcommand{\Subsection}{\subsection}
\providecommand{\nobreakdash}{\nobreak\hbox{-}}
\setlength{\emergencystretch}{2em}
\multlinegap0pt
\usepackage{amssymb}
\newtheorem{theorem}{Theorem}
\title{Group rings and hyperbolic geometry}
\author{Grigori Avramidi, Thomas Delzant}
\date{Source: arXiv:2309.16791v2 (CC BY 4.0)}
\begin{document}
\maketitle

\noindent\emph{Source and licence.} Group rings and hyperbolic geometry. Version arXiv:2309.16791v2 (CC BY 4.0), distributed by arXiv under the Creative Commons Attribution 4.0 International licence, http://creativecommons.org/licenses/by/4.0/, as recorded in the arXiv licence metadata for this exact version and shown on the abstract page https://arxiv.org/abs/2309.16791v2. Full licence notice: \texttt{licenses/arxiv-2309.16791-CC-BY-4.0.txt}.\par\smallskip

\noindent\textit{Editorial scope.} Counted unit: the article's theorem essential (source0.tex lines 1605--1608). The article's complete original proof of it is delivered with it (source0.tex lines 1609--1660, one \texttt{proof} environment of 52 lines). No mathematics is written, repaired or re-derived: the statement and the proof are the article's own lines, in the article's own notation, and no retained line is substituted or re-flowed.  Retained uncounted as the source-local context the proof needs: lines 1261--1270.  Nothing was omitted from any retained span, so the package carries no omission record.  No retained line contains a citation, so the wrapper carries no bibliography and no bibliography entry.  No retained line contains a figure environment, an image-inclusion command or drawing code, so no image is delivered and the package declares no asset dependency.  Editorial formatting only, in addition: an AMS \texttt{amsart} preamble that restates the article's own declarations and macros used by the retained lines, namely the environments \texttt{theorem} on separate counters, together with the standard packages those lines need.  The retained environments print in this document as Theorem 1 in order of appearance, whereas the article prints the same items with the numbering of its own full document.  The complete native source of the article as distributed in the arXiv e-print for this exact version is bundled unchanged as \texttt{sources/147723/source0.tex}.
\begin{theorem}
\label{freemodules}
Assume the group $G$ satisfies ${\mathcal{H}_{n}}$.
\begin{enumerate}
\item[1.] 
Every $n$-generated ideal in $\mathbb{K} [G]$ is a free $\mathbb{K}[G]$-module.
\item[2.] 
Every $n$-generated submodule of a free $\mathbb{K} [G]$-module is a free $\mathbb{K} [G]$-module.
\end{enumerate}
\end{theorem}

\begin{theorem}
\label{essential}
Assume that the group $G$ satisfies $\mathcal H_{n}$. If $X\rightarrow BG$ is a $d$-essential map, then $X$ has more than $n$ cells in each dimension $0<k<d$. 
\end{theorem}

\begin{proof}
Suppose not. Let $\hat X$ be the $G$-cover of $X$ induced by the map $f:X\rightarrow BG$. Then $\partial C_k(\hat X;\mathbb K)$ is an $n$-generated submodule of a free module, so it is a free module by Theorem \ref{freemodules}.2. Next, let $\iota$ denote the inclusion $\partial C_k(\hat X;\mathbb K)\hookrightarrow C_{k-1}(\hat X;\mathbb K)$. The image of $f_{k-1}\circ\iota:\partial C_k(\hat X;\mathbb K)\rightarrow C_{k-1}(EG;\mathbb K)$ is in the image of $\partial:C_k(EG;\mathbb K)\rightarrow C_{k-1}(EG;\mathbb K)$ since $f_*$ is a chain map $(f_{k-1}\partial=\partial f_k)$, so by lifting a free basis of $\partial C_k(\hat X;\mathbb K)$, we can construct a map $\lambda$ which makes the following diagram commute
%be the image of the map $f_{k-1}\circ\partial:C_k(\hat X;\mathbb K)\rightarrow C_{k-1}(EG;\mathbb K)$. Note that
%\begin{itemize}
%\item 
%$F$ is an $n$-generated submodule of a free module, hence it is a free module by Theorem \ref{freemodules}.2, and
%\item 
%$F$ lies in the image of the boundary map $\partial:C_k(EG;\mathbb K)\rightarrow C_{k-1}(EG;\mathbb K)$, since $f_*$ is a chain map ($\partial\circ f_k=f_{k-1}\circ\partial$).
%\end{itemize}
%Therefore, the inclusion $j:F\hookrightarrow C_{k-1}(EG;\mathbb K)$ has a lift $j':F\hookrightarrow C_k(EG;\mathbb K)$ such that $\partial\circ j'=j$. Next, let $i$ denote the inclusion $\partial C_k(\hat X;\mathbb K)\hookrightarrow C_{k-1}(\hat X;\mathbb K)$. Then the map $f_{k-1}\circ i:\partial C_k(\hat X;\mathbb K)\rightarrow C_{k-1}(EG;\mathbb K)$ lifts to $f':=j'\circ f_{k-1}\circ i:\partial C_k(\hat X;\mathbb K)\rightarrow C_k(EG;\mathbb K)$ satisfying $f_{k-1}\circ i =\partial\circ f'$. In summary, we have the following commutative diagram 
$$
\begin{array}{cccccccc}
\dots\rightarrow & C_{k+1}(\hat X;\mathbb K)&\rightarrow &C_k(\hat X;\mathbb K)&\stackrel{\partial}\rightarrow &C_{k-1}(\hat X;\mathbb K)&\rightarrow\dots\rightarrow &C_0(\hat X;\mathbb K)\\
&\downarrow&&\downarrow\partial&&||\hspace{0.5cm}&&||\\
\dots\rightarrow & 0&\rightarrow &\partial C_k(\hat X;\mathbb K)&\stackrel{\iota}\hookrightarrow &C_{k-1}(\hat X;\mathbb K)&\rightarrow\dots\rightarrow &C_0(\hat X;\mathbb K)\\
&\downarrow&&\downarrow \lambda&&\downarrow f_{k-1}&&\downarrow\\
\dots\rightarrow & C_{k+1}(EG;\mathbb K)&\rightarrow &C_k(EG;\mathbb K)&\stackrel{\partial}\rightarrow &C_{k-1}(EG;\mathbb K)&\rightarrow\dots\rightarrow &C_0(EG;\mathbb K).
\end{array}
$$
So, we have constructed a chain map $f'_*:C_*(\hat X;\mathbb K)\rightarrow C_*(EG;\mathbb K)$ which agrees with $f$ on $C_{<k}$ and is zero on $C_{>k}$, In particular, since $d>k$, the chain map $f'_*$ induces the zero map on $H^d(-;V)$. % factors as
%$$
%H^{d}(BG;V)\rightarrow H^{d}(T_*;V)\rightarrow H^{d}(X;V).
%$$
%The middle term is zero, so this map is the zero map. 

On the other hand, we claim that $f'_*$ is chain homotopic to $f_*$.
The chain homotopy $h_*:C_{*}(\hat X;\mathbb K)\rightarrow C_{*+1}(EG;\mathbb K)$ is constructed as follows. 

For every $i<k$, we set $h_i=0$. 

In order to define $h_k$ we remark that the image of $f_k-f_k'$ is in the kernel of $\partial$ since $f_j=f'_j$  for $j<k$. Since $EG$ is contractible, this implies that $f_k-f_k'$ maps to the image of $\partial$. As $C_k(\hat X;
\mathbb K)$ is a free module, $f_k-f_k'$ lifts to a map $h_k:C_k(\hat X;\mathbb K)\rightarrow C_{k+1}(EG;\mathbb K)$ such that $\partial h_k=f_k-f_k'$. 

In order to define $h_{k+1}$, we observe that $\partial(f_{k+1}-h_k\partial)=0$. Indeed, 
\begin{eqnarray*}
\partial(f_{k+1}-h_k\partial)&=&\partial f_{k+1}-\partial h_k\partial \\
&=&f_k\partial-(f_k-f_k')\partial\\
&=&f_k'\partial\\
&=&\partial f'_{k+1}=0,
\end{eqnarray*}
where we have used the fact that $f_*$ and $f'_*$ are chain maps, and that $f'_*$ is zero on $C_{>k}$. Therefore, since $EG$ is contractible, $f_{k+1}-h_k\partial$ maps to the image of $\partial$. So, it has a lift $h_{k+1}:C_{k+1}(\hat X;\mathbb K)\rightarrow C_{k+2}(EG;\mathbb K)$. By construction $f_{k+1}=h_k\partial+\partial h_{k+1}$. 

In higher degrees we argue by induction, assume that for $i>k$ the maps $h_{\leq i}$ have been defined, and compute
\begin{eqnarray*}
\partial(f_{i+1}-h_i\partial)&=&\partial f_{i+1}-\partial h_i\partial\\
&=&f_i\partial-(f_i-h_{i-1}\partial)\partial\\
&=&0.
\end{eqnarray*}
So we can define $h_{i+1}$ as a lift of $f_{i+1}-h_i\partial$. Proceeding in this way, we obtain $h_*$ satisfying $f_i-f'_i=\partial h_i+h_{i-1}\partial$ in for all $i$. Therefore $f_*$ is chain homotopic to $f'_*$.

We conclude that $f_*$ induces the zero map on $H^d(-;V),$ contradicting the hypothesis that $f:X\rightarrow BG$ is $d$-essential. 
\end{proof}

\end{document}
